\chapter{Nhật ký công việc hàng tuần}
\label{app:nhat-ky}

Mỗi dòng phải khớp với commit trong repo. Cột ``Người làm'' lấy theo tác giả commit.

\begin{table}[H]
  \centering
  \small
  \begin{tabularx}{\textwidth}{p{1.7cm} >{\hsize=1.1\hsize}X p{1.6cm} >{\hsize=0.9\hsize}X}
    \toprule
    \textbf{Tuần} & \textbf{Công việc đã làm} & \textbf{Người làm} & \textbf{Khó khăn / kết quả} \\
    \midrule
    1--2 & Lập nhóm, chọn công nghệ; cài môi trường, Hello World & --- & Chưa có commit trong repo \\
    3 (21--27/09) & Khởi tạo repo, dữ liệu Quy chế 2021; \texttt{node\_parser.py} tách theo Điều; sửa \texttt{check\_model\_memory.py}; bộ câu hỏi 01 và 02; dữ liệu sổ tay sinh viên; nháp báo cáo (PR \#1--\#6) & Phan Hữu Kha, Võ Thị Như Ý, Trần Hà Thanh Thùy, Huỳnh Thanh Tiến, solozy & Pipeline một văn bản chạy được \\
    4 (28/09--04/10) & Cấu hình 2 profile, RAG cơ bản, tiền xử lý Markdown hai tầng, \texttt{processed\_loader.py}, fingerprint chỉ mục, model server; báo cáo tuần 4 & solozy & 6 văn bản, 54 Node \\
    5 \newline (05--09/10) &
      Chạy LLM + model server trên Colab; hybrid BM25, baseline, kiểm định thống kê; 26 kiểm thử + CI;
      Gradio có nhật ký câu hỏi; chạy lại 56 câu; chú giải thang điểm trong prompt, \texttt{RAG\_THINKING};
      cập nhật Chương 3--5 &
      solozy &
      Truy hồi R@1 81\% / 90\%; rerank, hybrid không giúp; ``không tìm thấy'' 52\% $\to$ 41\%;
      phát hiện lỗi nhầm hệ 10/hệ 4 và lỗi không phân biệt khóa tuyển sinh \\
    \midrule
    6 \newline (10--16/10) &
      Hoàn thành gói nhiệm vụ N2 (D1: metadata hiệu lực \& manifest; D2: nạp dữ liệu đợt 3 gồm phúc khảo, tiết học, chuẩn đầu ra; D3: trích xuất 5 file CSV bảng quy đổi tiếng Anh; D5: bộ 15 câu hỏi kiểm thử dev đợt 3; R6: thay thế node Điều 9; R8: minh chứng Ollama benchmark và OOD; R12: phụ lục cá nhân) (PR \#3, branch \texttt{khaph}) &
      Phan Hữu Kha \newline (MSSV: 3123560038) &
      Corpus tăng lên 20 văn bản; chuẩn hóa 301 dòng quy đổi chứng chỉ; bộ test bao quát xung đột quy chế; hoàn thành 7 commit độc lập chuẩn conventional commits \\
    \bottomrule
  \end{tabularx}
  \caption{Nhật ký công việc theo tuần, lập từ \texttt{git log}.}
  \label{tab:nhat-ky-tuan}
\end{table}

\section{Tuần 5 (05--09/10/2026)}

\subsection{Công việc đã làm}

\begin{itemize}
  \item \textbf{Hạ tầng}: \texttt{scripts/setup\_server\_v2.sh}, \texttt{scripts/colab\_server.py} --- chạy
        Ollama và model server trên Colab, mở URL công khai qua Cloudflare tunnel, in sẵn các dòng
        \texttt{.env} cho máy local.
  \item \textbf{Truy hồi}: hybrid BM25 + vector (\texttt{src/bm25.py}, \texttt{QueryFusionRetriever}); đo
        bốn cấu hình và biến thể trọng số 3:1; baseline BM25 / không xếp hạng; khoảng tin cậy bootstrap và
        kiểm định McNemar (\texttt{eval/metrics.py}).
  \item \textbf{Kiểm thử}: 26 kiểm thử offline trong \texttt{tests/}, CI GitHub Actions, \texttt{Makefile}.
  \item \textbf{Giao diện}: Gradio (\texttt{app/app.py}) ghi mỗi câu hỏi vào
        \texttt{eval/results/questions\_log.jsonl}.
  \item \textbf{Đánh giá câu trả lời}: chạy lại 56 câu bằng \texttt{eval/run\_batch\_cau\_hoi.py} với prompt
        mới, chấm tay (Mục~\ref{sec:chay-lai-56-cau}).
  \item \textbf{Tầng sinh}: thêm chú giải thang điểm hệ 10 / hệ 4 vào \texttt{QA\_PROMPT}; thêm công tắc
        \texttt{RAG\_THINKING} (Mục~\ref{sec:nham-thang-diem}).
  \item \textbf{Báo cáo}: cập nhật Chương 3--5 và phụ lục này.
\end{itemize}

\subsection{Kết quả và khó khăn}

\begin{itemize}
  \item Truy hồi (vector): bộ 01 R@1 81\%, R@3 100\%; bộ 02 R@1 90\%. Đo lại ngày 08/10 cho cùng chỉ số.
  \item Rerank và hybrid không cải thiện trên 54 Node, nên để mặc định tắt (Chương~\ref{ch:nang-cao}).
  \item Phiên chạy lại 56 câu: tỷ lệ ``không tìm thấy'' giảm từ 52\% xuống 41\%, nhưng vẫn sai ở số tuần
        của học kỳ chính, điểm F và câu hỏi về chứng chỉ tiếng Anh; một câu trả về chuỗi \texttt{/no\_think}.
  \item Lỗi mới phát hiện: nhầm thang điểm 10 với thang điểm 4 (``A được quy đổi thành 4,0 điểm hệ 10''),
        và câu trả lời về điểm không phân biệt khóa tuyển sinh trước 2025 với từ 2025 trở đi (QĐ 3183).
  \item Khi đối chiếu lại khoản 4 Điều 9, nhóm phát hiện đáp án tham chiếu cũ của câu ``học bị F'' bị
        nhầm; đã sửa trong Bảng~\ref{tab:tra-loi}.
\end{itemize}

\subsection{Việc chuyển sang tuần 6}

\begin{enumerate}
  \item Chạy \texttt{make eval-answers} trên LLM Colab; chạy lại 56 câu theo từng yếu tố (prompt, top-$k$
        đưa vào LLM, \texttt{RAG\_THINKING}) để đo tác dụng của chú giải thang điểm.
  \item Thêm metadata khóa áp dụng cho mục sửa đổi của QĐ 3183 và yêu cầu câu trả lời nêu khóa.
  \item Lọc chuỗi \texttt{/no\_think} khỏi đầu ra và đếm như một lỗi khi chấm.
  \item Nạp văn bản quy đổi chứng chỉ tiếng Anh (thư mục \texttt{2\_quy\_doi\_tieng\_anh}): câu ``quy định
        về tiếng anh đầu ra'' trong nhật ký Gradio hiện trả lời ``không tìm thấy''.
\end{enumerate}

\section{Tuần 6 (10--16/10/2026) --- Phần việc Người 2 (Phan Hữu Kha)}

\subsection{Công việc đã hoàn thành}

Trong tuần 6, thành viên Phan Hữu Kha (MSSV: 3123560038, tài khoản GitHub: \texttt{phwkha}, nhánh \texttt{khaph}, PR \#3) đã hoàn thành toàn diện gói nhiệm vụ được phân công (D1, D2, D3, D5, R6, R8, R12):

\begin{itemize}
  \item \textbf{Nhiệm vụ D1 (Manifest hiệu lực \& Metadata phiên bản)}: Rà soát toàn bộ văn bản trong \texttt{data/processed/}, bổ sung các trường \texttt{ap\_dung\_tu\_khoa}, \texttt{ap\_dung\_den\_khoa}, \texttt{tinh\_trang\_hieu\_luc}, \texttt{ngay\_ban\_hanh}; tạo manifest đồng bộ \texttt{data/manifest.json} và tài liệu giải trình \texttt{docs/manifest\_hieu\_luc.md}.
  \item \textbf{Nhiệm vụ D2 (Nạp dữ liệu đợt 3)}: Số hóa và nạp 3 văn bản mới: Quy trình phúc khảo bài thi (định dạng F6), Khung giờ tiết học - phòng học (định dạng F4), và Chuẩn đầu ra tốt nghiệp năm 2020 (định dạng F4). Cập nhật biến \texttt{SGU\_HIEU\_LUC} trong \texttt{src/corpus.py}, nâng quy mô corpus từ 17 lên 20 văn bản.
  \item \textbf{Nhiệm vụ D3 (Trích xuất CSV các bảng quy đổi chứng chỉ tiếng Anh)}: Chuẩn hóa 5 văn bản quy đổi (2366, 2451, 2626, 2101, 2567) thành 5 file CSV trong \texttt{data/bang\_quy\_doi/} với tổng cộng 301 dòng dữ liệu, lược đồ chuẩn 8 cột, kiểm định số liệu khớp 100\% với văn bản gốc.
  \item \textbf{Nhiệm vụ D5 (Bộ câu hỏi kiểm thử dev đợt 3)}: Xây dựng \texttt{eval/dev_questions_dot3.json} gồm 15 câu hỏi bao quát 4 chủ đề trọng điểm (phúc khảo, tiết học, chuẩn đầu ra, kiểm thử xung đột phiên bản quy chế) và tích hợp vào pipeline đánh giá \texttt{eval/evaluate.py}.
  \item \textbf{Nhiệm vụ R6 (Báo cáo - Node Điều 9)}: Thay thế placeholder \texttt{<nguyên văn Điều 9>} tại Chương 5 (\texttt{report/chapters/05-node-parser.tex}) bằng nội dung trích xuất thực tế từ Quy chế 2021.
  \item \textbf{Nhiệm vụ R8 (Báo cáo - Minh chứng thực nghiệm mô hình)}: Bổ sung số liệu thực nghiệm đo đạc độ trễ và thông lượng trên Ollama (Qwen3:8b), phân tích cơ chế từ chối câu hỏi ngoài miền (OOD) và giải pháp xử lý triệt để lỗi rò rỉ chỉ thị \texttt{/no\_think} (Chương 7 và Chương 10).
  \item \textbf{Nhiệm vụ R12 (Phụ lục và tổng kết cá nhân)}: Hoàn thiện bảng nhật ký công việc cá nhân, phương hướng nghiên cứu phát triển (Chương 12) và bản công bố sử dụng AI (Chương 11).
\end{itemize}

\subsection{Danh mục Commit trên nhánh \texttt{khaph}}

Toàn bộ các nhiệm vụ trên được đóng gói thành các commit độc lập, tuân thủ nghiêm ngặt quy chuẩn Conventional Commits:
\begin{enumerate}
  \item \texttt{50e73f0}: \texttt{feat(data): hoan thanh D1 - xac minh hieu luc va cap nhat metadata phien ban}
  \item \texttt{bfe96f9}: \texttt{feat(data): hoan thanh D2 - nap du lieu dot 3 (phuc khao, tiet hoc, chuan dau ra)}
  \item \texttt{7bf3a27}: \texttt{feat(data): hoan thanh D3 - trich xuat CSV cac bang quy doi chung chi tieng Anh}
  \item \texttt{a54ea28}: \texttt{feat(eval): hoan thanh D5 - bo cau hoi kiem thu dev dot 3}
  \item \texttt{5e3bbf9}: \texttt{docs(report): hoan thanh R6 - thay placeholder Dieu 9 bang Node thuc te}
  \item \texttt{68cdcf0}: \texttt{docs(report): hoan thanh R8 - bo sung minh chung thuc nghiem mo hinh ngon ngu}
  \item \texttt{docs(report): hoan thanh R12 - nhat ky ca nhan va phu luc A cua N2}
\end{enumerate}
